\documentclass[10pt,a4paper]{article}
\usepackage[utf8]{inputenc}
\usepackage[T1]{fontenc}
\usepackage[brazil]{babel}
\usepackage{geometry}
\usepackage{enumitem}
\usepackage{parskip}
\geometry{margin=1.5cm}
\setlength{\parindent}{0pt}

\begin{document}

\begin{center}
    {\Large \textbf{Atividade Prática A1 -- Grafos}}\\[2mm]
    INE5413 -- Ciências da Computação -- UFSC\\
    \textbf{Integrantes:} Nome 1, Nome 2, Nome 3
\end{center}

\textbf{Item 1 -- Representação do grafo.} Foi implementada a classe \texttt{Graph} para representar um grafo não-dirigido e ponderado. A estrutura principal adotada foi uma lista 1-indexada de dicionários de adjacência, na qual \texttt{adj[u][v]} armazena diretamente o peso da aresta \(\{u,v\}\). Essa escolha permite tempo \(O(1)\) médio para as operações \texttt{haAresta(u,v)}, \texttt{peso(u,v)}, \texttt{grau(v)} e \texttt{rotulo(v)}, atendendo ao requisito do enunciado sempre que possível. Os rótulos dos vértices foram armazenados em uma lista 1-indexada separada, também com acesso \(O(1)\). A leitura do arquivo é linear no tamanho da entrada, isto é, \(O(|V|+|E|)\).

\textbf{Item 2 -- Busca em largura (BFS).} Para a BFS foram utilizados um vetor booleano de visitados, um vetor de níveis e uma fila \texttt{deque}. A fila garante inserção e remoção no início/fim em \(O(1)\), tornando a busca global \(O(|V|+|E|)\). Como o exercício exige a saída agrupada por nível, os vértices descobertos foram adicionados em uma lista de listas, preservando uma organização simples para impressão posterior. Os pesos do arquivo são ignorados, conforme solicitado.

\textbf{Item 3 -- Ciclo euleriano.} A solução primeiro verifica duas condições necessárias e suficientes para grafos não-dirigidos: todos os vértices de grau não nulo devem pertencer ao mesmo componente conexo e todos os graus devem ser pares. Caso essas condições sejam satisfeitas, o ciclo é construído com o algoritmo de Hierholzer, que percorre cada aresta essencialmente uma vez, com custo \(O(|V|+|E|)\). Foi utilizada uma estrutura auxiliar de multiplicidade de arestas para consumir as arestas durante a construção do ciclo sem alterar a representação original do grafo.

\textbf{Item 4 -- Menores caminhos de fonte única.} Foi escolhido o algoritmo de Dijkstra, implementado com fila de prioridade baseada em \texttt{heapq}. Essa estratégia é adequada para pesos não negativos e apresenta complexidade \(O((|V|+|E|)\log |V|)\). As distâncias mínimas são guardadas em um vetor \texttt{dist} e os predecessores em um vetor \texttt{parent}, permitindo reconstruir o caminho mínimo até cada destino exigido pelo enunciado. A impressão final mostra, para cada vértice, o caminho completo a partir da origem e a respectiva distância acumulada.

\textbf{Item 5 -- Floyd-Warshall.} Para os menores caminhos entre todos os pares, foi adotada uma matriz de distâncias \(dist[i][j]\), inicializada com infinito, zero na diagonal principal e peso das arestas nas posições correspondentes. O algoritmo de Floyd-Warshall foi então aplicado com três laços aninhados, resultando em complexidade \(O(|V|^3)\) e uso de memória \(O(|V|^2)\). Apesar de mais custoso que Dijkstra para uma única origem, ele é apropriado quando se deseja obter todas as distâncias entre todos os pares de vértices em uma única execução.

\textbf{Validação.} Antes da execução em ambiente fechado, foram construídos testes unitários com \texttt{unittest} para validar a leitura do grafo, a BFS, a detecção/construção de ciclo euleriano, o Dijkstra e o Floyd-Warshall. Dessa forma, a solução foi verificada automaticamente em cenários pequenos e controlados, reduzindo a chance de erro funcional na entrega.

\end{document}
